% ══════════════════════════════════════════════════════════════════════════════
%                         فصل سوم
%                 اقتصاد سیاسی و جدال بنیادین
% ══════════════════════════════════════════════════════════════════════════════

\chapter{اقتصاد سیاسی و جدال بنیادین}
\label{ch:political-economy}

\begin{flushright}
\textit{«اقتصاد، ادامه سیاست با ابزارهای دیگر است.»}
\end{flushright}

% ══════════════════════════════════════════════════════════════════════════════
%                         مقدمه فصل
% ══════════════════════════════════════════════════════════════════════════════

\lettrine[lines=3, lraise=0.1, nindent=0.5em]{\textcolor{CenterGreen}{ا}}{گر} 
در فصل پیش دیدیم که چپ و راست ریشه‌های فلسفی متفاوتی دارند، در این فصل به میدان اصلی نبرد می‌رسیم: \keyword{اقتصاد}. پرسش‌هایی مانند «مالکیت از آنِ کیست؟»، «بازار آزاد باشد یا کنترل‌شده؟» و «دولت چه نقشی در اقتصاد دارد؟» از قرن نوزدهم به این‌سو، محور اصلی اختلاف چپ و راست بوده‌اند. درک این جدال‌ها، کلید فهم سیاست مدرن است.

% ══════════════════════════════════════════════════════════════════════════════
\section{مالکیت: بنیادی‌ترین اختلاف}
% ══════════════════════════════════════════════════════════════════════════════

هیچ مسئله‌ای به اندازه \keyword{مالکیت} چپ و راست را از هم جدا نکرده است. آیا مالکیت خصوصی حق طبیعی است یا دزدی؟

\subsection{دفاع از مالکیت خصوصی}

سنت لیبرالی و محافظه‌کار، مالکیت خصوصی را بنیان جامعه آزاد می‌داند.

\begin{personbox}[جان لاک و نظریه کار]

\person{جان لاک} اولین نظریه مدرن مالکیت را ارائه داد:

\begin{enumerate}[nosep]
    \item خداوند زمین را به بشر مشترکاً داده است
    \item اما هر کس مالک بدن و کار خودش هست
    \item وقتی کار خود را با طبیعت می‌آمیزی، آن را از حالت مشترک خارج می‌کنی
    \item پس مالکیت خصوصی از \textbf{کار} ناشی می‌شود
\end{enumerate}

\textbf{شرط لاک:} تا زمانی که «به اندازه کافی و به همان خوبی» برای دیگران باقی بماند.

\mediumspace

\textbf{نقد مارکسی:} اما وقتی زمین تمام شد چه؟ شرط لاک دیگر برقرار نیست.
\end{personbox}

\begin{quotebox}[آدام اسمیت]
\textit{«هیچ‌چیز به اندازه مالکیت امن، تشویق به کار و صرفه‌جویی نمی‌کند. هر انسان طبیعتاً تلاش می‌کند وضعیت خود را بهبود بخشد.»}
\end{quotebox}

\begin{defbox}[استدلال‌های له مالکیت خصوصی]
\begin{enumerate}[nosep]
    \item \textbf{استدلال حق طبیعی (لاک):} مالکیت از کار ناشی می‌شود
    \item \textbf{استدلال کارایی (اسمیت):} انگیزه کار و نوآوری ایجاد می‌کند
    \item \textbf{استدلال آزادی (هایک):} پایه استقلال فردی در برابر دولت است
    \item \textbf{استدلال محافظه‌کارانه (برک):} نهاد تاریخی محترم است
\end{enumerate}
\end{defbox}

\subsection{نقد مالکیت خصوصی}

در مقابل، سنت سوسیالیستی و آنارشیستی مالکیت خصوصیِ ابزار تولید را به چالش کشیده است.

\begin{personbox}[پیر ژوزف پرودون (۱۸۰۹-۱۸۶۵)]

\textbf{ملیت:} فرانسوی

\textbf{اثر کلیدی:} \book{مالکیت چیست؟} (۱۸۴۰)

\textbf{جمله معروف:} «مالکیت دزدی است!»

\mediumspace

\person{پرودون} اولین کسی بود که خود را «آنارشیست» نامید. او استدلال کرد:
\begin{itemize}[nosep]
    \item مالکیت بزرگ (زمین، کارخانه) از استثمار و زور ناشی شده
    \item مالک، بدون کار، از کار دیگران سود می‌برد (اجاره، سود)
    \item این «دزدی» قانونی شده است
\end{itemize}

\mediumspace

\textbf{راه‌حل پرودون:} نه مالکیت خصوصی بزرگ، نه مالکیت دولتی، بلکه «تصرف» و تعاونی‌ها.
\end{personbox}

\begin{personbox}[کارل مارکس (۱۸۱۸-۱۸۸۳)]

\textbf{ملیت:} آلمانی

\textbf{آثار کلیدی:} \book{سرمایه} (۱۸۶۷)، \book{مانیفست کمونیست} (۱۸۴۸)

\textbf{ایده محوری:} نقد اقتصاد سیاسی سرمایه‌داری

\mediumspace

\person{مارکس} نقد پرودون را علمی‌تر کرد:
\begin{itemize}[nosep]
    \item مالکیت ابزار تولید، طبقه سرمایه‌دار را از طبقه کارگر جدا می‌کند
    \item کارگر نیروی کار می‌فروشد، اما ارزشی بیش از دستمزدش تولید می‌کند
    \item این \term{ارزش اضافی} به جیب سرمایه‌دار می‌رود = \textbf{استثمار}
    \item مالکیت اشتراکی ابزار تولید، استثمار را پایان می‌دهد
\end{itemize}

\mediumspace

\textbf{نقل‌قول:} \textit{«سرمایه، کار مُرده است که مانند خون‌آشام، فقط با مکیدن کار زنده زنده می‌ماند.»}
\end{personbox}

\begin{figure}[htbp]
\centering
\begin{tikzpicture}[scale=0.85]

% عنوان
\node[font=\large\bfseries, text=DarkGray] at (0, 7) {\fa{طیف نظرات درباره مالکیت}};

% محور
\draw[line width=4pt, MediumGray] (-7, 4.5) -- (7, 4.5);

% نقاط روی طیف
% مالکیت مطلق
\fill[RightBlueDark] (6, 4.5) circle (10pt);
\node[below, text width=2.5cm, align=center, text=DarkGray, font=\small] at (6, 3.7) {
    \fa{مالکیت مطلق}\\
    \fa{(لیبرتارین)}
};

% مالکیت محدود
\fill[RightBlueLight] (3, 4.5) circle (10pt);
\node[below, text width=2.5cm, align=center, text=DarkGray, font=\small] at (3, 3.7) {
    \fa{مالکیت محدود}\\
    \fa{(لیبرال/محافظه‌کار)}
};

% مالکیت مشروط
\fill[CenterGreen] (0, 4.5) circle (10pt);
\node[below, text width=2.5cm, align=center, text=DarkGray, font=\small] at (0, 3.7) {
    \fa{مالکیت مشروط}\\
    \fa{(سوسیال‌دموکرات)}
};

% مالکیت محدود ابزار تولید
\fill[LeftRedLight] (-3, 4.5) circle (10pt);
\node[below, text width=2.5cm, align=center, text=DarkGray, font=\small] at (-3, 3.7) {
    \fa{ملی‌سازی بخشی}\\
    \fa{(سوسیالیست دموکراتیک)}
};

% الغای مالکیت خصوصی
\fill[LeftRedDark] (-6, 4.5) circle (10pt);
\node[below, text width=2.5cm, align=center, text=DarkGray, font=\small] at (-6, 3.7) {
    \fa{الغای مالکیت}\\
    \fa{(کمونیست)}
};

% برچسب‌ها
\node[text=RightBlue, font=\bfseries] at (7.5, 4.5) {\fa{راست}};
\node[text=LeftRed, font=\bfseries] at (-7.5, 4.5) {\fa{چپ}};

% توضیح پایین
\node[text=MediumGray, font=\footnotesize, text width=12cm, align=center] at (0, 2.2) {
    \fa{توجه: تمایز مهم بین «مالکیت شخصی» (خانه، اتومبیل) و «مالکیت ابزار تولید» (کارخانه، زمین کشاورزی) است.}\\
    \fa{اغلب سوسیالیست‌ها فقط دومی را هدف می‌گیرند.}
};

\end{tikzpicture}
\caption{طیف دیدگاه‌ها درباره مالکیت}
\label{fig:property-spectrum}
\end{figure}

\begin{debatebox}[مالکیت: حق یا دزدی؟]
\textbf{مدافعان مالکیت:}
\begin{itemize}[nosep]
    \item مالکیت از کار مشروع ناشی می‌شود
    \item بدون مالکیت امن، انگیزه تولید از بین می‌رود
    \item مالکیت خصوصی، آزادی فردی را تضمین می‌کند
    \item جوامع بدون مالکیت خصوصی، فقیر و استبدادی شدند
\end{itemize}

\textbf{منتقدان مالکیت (ابزار تولید):}
\begin{itemize}[nosep]
    \item ثروت اولیه اغلب از خشونت و استعمار آمده
    \item مالکیت بزرگ از کار دیگران تغذیه می‌شود
    \item «آزادی» قرارداد، وقتی یک طرف چیزی ندارد، بی‌معناست
    \item شوروی نماینده واقعی ایده نبود
\end{itemize}
\end{debatebox}

% ══════════════════════════════════════════════════════════════════════════════
\section{بازار در برابر برنامه‌ریزی}
% ══════════════════════════════════════════════════════════════════════════════

\subsection{بازار آزاد: دست نامرئی}

\begin{personbox}[آدام اسمیت (۱۷۲۳-۱۷۹۰)]

\textbf{ملیت:} اسکاتلندی

\textbf{اثر کلیدی:} \book{ثروت ملل} (۱۷۷۶)

\textbf{ایده محوری:} دست نامرئی بازار

\mediumspace

\person{آدام اسمیت}، پدر اقتصاد مدرن، استدلال کرد که وقتی هر فرد منفعت خود را دنبال می‌کند، «دست نامرئی» بازار، منفعت عمومی را تأمین می‌کند.

\textbf{نقل‌قول معروف:} \textit{«این از خیرخواهی قصاب، آبجوساز یا نانوا نیست که ما انتظار شام خود را داریم، بلکه از توجه آنان به منفعت خودشان.»}

\mediumspace

\textbf{نکته مهم:} اسمیت طرفدار «بازار افسارگسیخته» نبود. او از آموزش عمومی، زیرساخت‌ها و حتی محدودیت‌هایی بر بانک‌ها دفاع کرد.
\end{personbox}

\begin{defbox}[تعریف: مکانیسم بازار]
در نظریه کلاسیک:
\begin{itemize}[nosep]
    \item \textbf{قیمت‌ها} اطلاعات را منتقل می‌کنند
    \item \textbf{عرضه و تقاضا} قیمت‌ها را تعیین می‌کنند
    \item \textbf{سود و زیان} تولیدکنندگان را راهنمایی می‌کند
    \item \textbf{رقابت} کیفیت را بالا و قیمت را پایین می‌برد
    \item نیازی به برنامه‌ریز مرکزی نیست
\end{itemize}
\end{defbox}

\subsection{نقد بازار: شکست‌های بازار}

اقتصاددانان، حتی طرفداران بازار، می‌دانند که بازار همیشه کارآمد نیست.

\begin{table}[htbp]
\centering
\caption{انواع شکست بازار}
\label{tab:market-failures}
\begin{tabular}{r|r|r}
\toprule
\textbf{نوع شکست} & \textbf{توضیح} & \textbf{مثال} \\
\midrule
انحصار & یک فروشنده قیمت را دیکته می‌کند & شرکت‌های بزرگ فناوری \\
\midrule
برون‌دادهای منفی & هزینه به دیگران تحمیل می‌شود & آلودگی محیط زیست \\
\midrule
کالاهای عمومی & همه استفاده می‌کنند، کسی نمی‌پردازد & دفاع ملی، چراغ دریایی \\
\midrule
اطلاعات نامتقارن & یک طرف بیشتر می‌داند & بیمه، بازار دست‌دوم \\
\midrule
نابرابری & بازار عدالت را تضمین نمی‌کند & فقر در کنار ثروت \\
\bottomrule
\end{tabular}
\end{table}

\subsection{برنامه‌ریزی مرکزی: جایگزین سوسیالیستی}

برخی سوسیالیست‌ها پیشنهاد کردند به جای «آنارشی تولید» سرمایه‌داری، اقتصاد را به صورت \keyword{عقلانی و برنامه‌ریزی‌شده} اداره کنیم.

\begin{historybox}[اقتصاد برنامه‌ریزی‌شده در عمل]
اتحاد شوروی از ۱۹۲۸ تا ۱۹۹۱ اقتصاد برنامه‌ریزی مرکزی داشت:
\begin{itemize}[nosep]
    \item «گوسپلان» (کمیته برنامه‌ریزی) تولید را تعیین می‌کرد
    \item برنامه‌های پنج‌ساله اهداف را مشخص می‌کردند
    \item قیمت‌ها توسط دولت تعیین می‌شدند
\end{itemize}
\textbf{دستاوردها:} صنعتی‌شدن سریع، پیروزی در جنگ جهانی دوم، فضانوردی

\textbf{مشکلات:} کمبود مزمن کالاها، ناکارآمدی، فقدان نوآوری، فساد
\end{historybox}

\subsection{مناظره محاسبه سوسیالیستی}

یکی از مهم‌ترین مناظرات اقتصادی قرن بیستم، \term{مناظره محاسبه سوسیالیستی} بود.

\begin{figure}[htbp]
\centering
\begin{tikzpicture}[scale=0.85]

% عنوان
\node[font=\large\bfseries, text=DarkGray] at (0, 6.5) {\fa{مناظره محاسبه سوسیالیستی (دهه ۱۹۲۰-۱۹۴۰)}};

% طرف راست - میزس/هایک
\fill[RightBlue, opacity=0.2, rounded corners=10pt] (0.5, -0.5) rectangle (7, 5.5);
\draw[line width=2pt, color=RightBlue, rounded corners=10pt] (0.5, -0.5) rectangle (7, 5.5);

\node[text=RightBlue, font=\bfseries\large] at (3.75, 5) {\fa{میزس و هایک}};

\node[text=DarkGray, font=\small, text width=5.5cm, align=right] at (3.75, 3.5) {
    \fa{بدون قیمت‌های بازار،}\\
    \fa{محاسبه اقتصادی غیرممکن است.}\\[5pt]
    \fa{قیمت‌ها اطلاعات پراکنده را}\\
    \fa{جمع‌آوری می‌کنند.}\\[5pt]
    \fa{هیچ برنامه‌ریزی نمی‌تواند}\\
    \fa{جای میلیون‌ها تصمیم را بگیرد.}
};

\node[text=RightBlueDark, font=\small\bfseries] at (3.75, 0.3) {\fa{نتیجه: سوسیالیسم ناممکن است}};

% طرف چپ - لانگه/لرنر
\fill[LeftRed, opacity=0.2, rounded corners=10pt] (-7, -0.5) rectangle (-0.5, 5.5);
\draw[line width=2pt, color=LeftRed, rounded corners=10pt] (-7, -0.5) rectangle (-0.5, 5.5);

\node[text=LeftRed, font=\bfseries\large] at (-3.75, 5) {\fa{لانگه و لرنر}};

\node[text=DarkGray, font=\small, text width=5.5cm, align=right] at (-3.75, 3.5) {
    \fa{می‌توان «سوسیالیسم بازار»}\\
    \fa{طراحی کرد.}\\[5pt]
    \fa{دولت قیمت‌ها را آزمون‌وخطایی}\\
    \fa{تعیین می‌کند.}\\[5pt]
    \fa{مدیران دولتی مانند}\\
    \fa{مدیران خصوصی عمل کنند.}
};

\node[text=LeftRedDark, font=\small\bfseries] at (-3.75, 0.3) {\fa{نتیجه: سوسیالیسم ممکن است}};

% علامت مناظره
\node[text=GoldAccent, font=\Huge\bfseries] at (0, 2.5) {\fa{⚔}};

\end{tikzpicture}
\caption{مناظره محاسبه سوسیالیستی}
\label{fig:socialist-calculation}
\end{figure}

\begin{personbox}[فردریش هایک (۱۸۹۹-۱۹۹۲)]

\textbf{ملیت:} اتریشی-بریتانیایی

\textbf{آثار کلیدی:} \book{راه بردگی} (۱۹۴۴)، \book{استفاده از دانش در جامعه} (۱۹۴۵)

\textbf{جایزه نوبل اقتصاد:} ۱۹۷۴

\mediumspace

\person{هایک} استدلال میزس را عمیق‌تر کرد:
\begin{itemize}[nosep]
    \item دانش در جامعه \textbf{پراکنده} است—هیچ‌کس همه را نمی‌داند
    \item قیمت‌ها این دانش پراکنده را \textbf{تجمیع} می‌کنند
    \item هیچ برنامه‌ریز مرکزی نمی‌تواند این کار را انجام دهد
    \item برنامه‌ریزی مرکزی ضرورتاً به \textbf{استبداد} می‌انجامد
\end{itemize}

\mediumspace

\textbf{نقل‌قول:} \textit{«کنجکاو است که آدمی که از خودش شروع می‌کند، معتقد است می‌تواند دانشی کسب کند که او را قادر به شکل دادن به جامعه کند.»}
\end{personbox}

% ══════════════════════════════════════════════════════════════════════════════
\section{نقش دولت در اقتصاد}
% ══════════════════════════════════════════════════════════════════════════════

حتی اگر مالکیت خصوصی و بازار را بپذیریم، همچنان پرسش بزرگی باقی است: دولت چه نقشی باید در اقتصاد داشته باشد؟

\subsection{طیف نقش‌های دولت}

\begin{figure}[htbp]
\centering
\begin{tikzpicture}[scale=0.75]

% عنوان
\node[font=\large\bfseries, text=DarkGray] at (0, 8) {\fa{طیف نقش دولت در اقتصاد}};

% محور عمودی - میزان دخالت
\draw[line width=3pt, MediumGray, -{Stealth[length=10pt]}] (-7, 0) -- (-7, 7);
\node[rotate=90, text=MediumGray, font=\bfseries] at (-7.8, 3.5) {\fa{میزان دخالت دولت}};

% سطوح مختلف
% سطح ۱: دولت حداقلی
\fill[RightBlueDark, opacity=0.8, rounded corners=5pt] (-6.5, 0.5) rectangle (6.5, 1.5);
\node[text=white, font=\bfseries] at (0, 1) {\fa{دولت حداقلی (شب‌پاس)}};
\node[text=VeryLightGray, font=\scriptsize] at (0, 0.6) {\fa{فقط امنیت، دادگستری، دفاع}};

% سطح ۲: دولت لیبرال کلاسیک
\fill[RightBlueLight, opacity=0.8, rounded corners=5pt] (-6.5, 1.8) rectangle (6.5, 2.8);
\node[text=white, font=\bfseries] at (0, 2.5) {\fa{دولت لیبرال کلاسیک}};
\node[text=VeryLightGray, font=\scriptsize] at (0, 2.1) {\fa{+ زیرساخت، آموزش پایه، بانک مرکزی}};

% سطح ۳: دولت رفاه
\fill[CenterGreen, opacity=0.8, rounded corners=5pt] (-6.5, 3.1) rectangle (6.5, 4.1);
\node[text=white, font=\bfseries] at (0, 3.8) {\fa{دولت رفاه (کینزی)}};
\node[text=VeryLightGray, font=\scriptsize] at (0, 3.4) {\fa{+ تأمین اجتماعی، بهداشت، سیاست پولی و مالی فعال}};

% سطح ۴: دولت توسعه‌گرا
\fill[LeftRedLight, opacity=0.8, rounded corners=5pt] (-6.5, 4.4) rectangle (6.5, 5.4);
\node[text=white, font=\bfseries] at (0, 5.1) {\fa{دولت توسعه‌گرا}};
\node[text=VeryLightGray, font=\scriptsize] at (0, 4.7) {\fa{+ صنایع استراتژیک دولتی، برنامه‌ریزی هدفمند}};

% سطح ۵: دولت سوسیالیستی
\fill[LeftRedDark, opacity=0.8, rounded corners=5pt] (-6.5, 5.7) rectangle (6.5, 6.7);
\node[text=white, font=\bfseries] at (0, 6.4) {\fa{دولت سوسیالیستی}};
\node[text=VeryLightGray, font=\scriptsize] at (0, 6) {\fa{مالکیت عمومی ابزار تولید، برنامه‌ریزی مرکزی}};

% فلش‌های کناری
\node[text=RightBlue, font=\small, rotate=90] at (7.3, 1.2) {\fa{لیبرتارین}};
\node[text=CenterGreen, font=\small, rotate=90] at (7.3, 3.6) {\fa{سوسیال‌دموکرات}};
\node[text=LeftRed, font=\small, rotate=90] at (7.3, 6.2) {\fa{کمونیست}};

\end{tikzpicture}
\caption{سطوح مختلف دخالت دولت در اقتصاد}
\label{fig:state-intervention}
\end{figure}

\subsection{دولت رفاه: میانه کینزی}

\begin{personbox}[جان مینارد کینز (۱۸۸۳-۱۹۴۶)]

\textbf{ملیت:} بریتانیایی

\textbf{اثر کلیدی:} \book{نظریه عمومی اشتغال، بهره و پول} (۱۹۳۶)

\textbf{ایده محوری:} تقاضای مؤثر و نقش دولت در تثبیت اقتصاد

\mediumspace

\person{کینز} در بحران بزرگ ۱۹۲۹ نظریه‌ای انقلابی ارائه داد:
\begin{itemize}[nosep]
    \item بازار خودبه‌خود به تعادل اشتغال کامل نمی‌رسد
    \item ممکن است «تعادل» با بیکاری بالا شکل بگیرد
    \item در رکود، دولت باید با \textbf{هزینه‌کرد} تقاضا ایجاد کند
    \item در رونق، دولت باید با مالیات و کاهش هزینه جلو تورم را بگیرد
\end{itemize}

\mediumspace

\textbf{نقل‌قول:} \textit{«در بلندمدت همه ما مرده‌ایم.»}

\mediumspace

\textbf{پیامد:} توجیه فکری دولت رفاه و سیاست‌های تثبیتی
\end{personbox}

\begin{historybox}[عصر طلایی کینزگرایی (۱۹۴۵-۱۹۷۳)]
پس از جنگ جهانی دوم، اغلب کشورهای غربی سیاست‌های کینزی را پذیرفتند:
\begin{itemize}[nosep]
    \item دولت رفاه: بیمه بیکاری، بازنشستگی، بهداشت عمومی
    \item سیاست مالی فعال: هزینه‌کرد دولت برای تثبیت اقتصاد
    \item اقتصاد مختلط: بخش خصوصی + صنایع ملی‌شده
\end{itemize}
\textbf{نتیجه:} رشد اقتصادی بی‌سابقه، کاهش نابرابری، اشتغال کامل تقریبی

این دوره «سازش بزرگ» یا «اجماع پسا جنگ» نامیده می‌شود.
\end{historybox}

% ══════════════════════════════════════════════════════════════════════════════
\section{نظریه‌های ارزش: اختلافی بنیادین}
% ══════════════════════════════════════════════════════════════════════════════

پشت اختلافات سیاسی درباره اقتصاد، اختلافی نظری درباره «ارزش» نهفته است: ارزش کالاها از کجا می‌آید؟

\subsection{نظریه ارزش کار}

\begin{defbox}[تعریف: نظریه ارزش کار (LTV)]
این نظریه که از لاک، اسمیت و ریکاردو آمده و مارکس آن را کامل کرد، می‌گوید:
\begin{itemize}[nosep]
    \item ارزش یک کالا برابر است با مقدار \textbf{کار اجتماعاً لازم} برای تولید آن
    \item سرمایه‌دار بخشی از این ارزش (ارزش اضافی) را بدون کار تصاحب می‌کند
    \item این «استثمار» است—حتی اگر کارگر داوطلبانه قرارداد ببندد
\end{itemize}
\textbf{پیامد سیاسی:} سود سرمایه‌دار ناعادلانه است.
\end{defbox}

\subsection{نظریه ارزش ذهنی}

\begin{defbox}[تعریف: نظریه ارزش ذهنی (مارژینالیسم)]
انقلاب مارژینالیستی (دهه ۱۸۷۰) نظریه ارزش کار را رد کرد:
\begin{itemize}[nosep]
    \item ارزش \textbf{ذهنی} است—بستگی به ترجیحات افراد دارد
    \item ارزش = \textbf{مطلوبیت نهایی} (آخرین واحد چقدر ارزشمند است؟)
    \item قیمت‌ها از تعامل عرضه و تقاضا شکل می‌گیرند
    \item کار، هزینه تولید است، نه منبع ارزش
\end{itemize}
\textbf{پیامد سیاسی:} سود، پاداش ریسک و کارآفرینی است، نه استثمار.
\end{defbox}

\begin{comparebox}[مقایسه: دو نظریه ارزش]
\begin{table}[H]
\centering
\begin{tabular}{r|c|c}
\toprule
\textbf{موضوع} & \textbf{ارزش کار} & \textbf{ارزش ذهنی} \\
\midrule
منبع ارزش & کار بشری & ترجیحات فردی \\
قیمت & بازتاب ارزش واقعی & تعادل عرضه و تقاضا \\
سود سرمایه & استثمار & پاداش ریسک و انتظار \\
دستمزد & کمتر از ارزش تولید & قیمت نیروی کار \\
پیامد سیاسی & نقد سرمایه‌داری & توجیه سرمایه‌داری \\
طرفداران & مارکسیست‌ها & نئوکلاسیک‌ها، اتریشی‌ها \\
\bottomrule
\end{tabular}
\end{table}
\end{comparebox}

% ══════════════════════════════════════════════════════════════════════════════
\section{مکاتب اقتصادی معاصر}
% ══════════════════════════════════════════════════════════════════════════════

در قرن بیستم، مکاتب اقتصادی متعددی شکل گرفتند که هر کدام رابطه‌ای با طیف سیاسی دارند.

\begin{figure}[htbp]
\centering
\begin{tikzpicture}[scale=0.8]

% عنوان
\node[font=\large\bfseries, text=DarkGray] at (0, 8.5) {\fa{تایم‌لاین: تحول اندیشه اقتصادی}};

% خط زمان
\draw[line width=4pt, MediumGray] (-7, 6) -- (7, 6);

% نقاط کلیدی
% اقتصاد کلاسیک
\fill[MediumGray] (-6, 6) circle (8pt);
\node[above, text=DarkGray, font=\small\bfseries] at (-6, 6.3) {\fa{۱۷۷۶}};
\node[below, text width=2cm, align=center, text=DarkGray, font=\scriptsize] at (-6, 5.4) {
    \fa{اسمیت}\\
    \fa{کلاسیک}
};

% مارکس
\fill[LeftRed] (-4, 6) circle (8pt);
\node[above, text=LeftRed, font=\small\bfseries] at (-4, 6.3) {\fa{۱۸۶۷}};
\node[below, text width=2cm, align=center, text=DarkGray, font=\scriptsize] at (-4, 5.4) {
    \fa{مارکس}\\
    \fa{سرمایه}
};

% مارژینالیسم
\fill[RightBlueLight] (-2, 6) circle (8pt);
\node[above, text=RightBlueLight, font=\small\bfseries] at (-2, 6.3) {\fa{۱۸۷۰s}};
\node[below, text width=2cm, align=center, text=DarkGray, font=\scriptsize] at (-2, 5.4) {
    \fa{انقلاب}\\
    \fa{مارژینالیستی}
};

% کینز
\fill[CenterGreen] (0, 6) circle (8pt);
\node[above, text=CenterGreen, font=\small\bfseries] at (0, 6.3) {\fa{۱۹۳۶}};
\node[below, text width=2cm, align=center, text=DarkGray, font=\scriptsize] at (0, 5.4) {
    \fa{کینز}\\
    \fa{نظریه عمومی}
};

% هایک/فریدمن
\fill[RightBlue] (2.5, 6) circle (8pt);
\node[above, text=RightBlue, font=\small\bfseries] at (2.5, 6.3) {\fa{۱۹۴۴-۶۲}};
\node[below, text width=2cm, align=center, text=DarkGray, font=\scriptsize] at (2.5, 5.4) {
    \fa{هایک، فریدمن}\\
    \fa{ضد کینز}
};

% نئولیبرالیسم
\fill[RightBlueDark] (5, 6) circle (8pt);
\node[above, text=RightBlueDark, font=\small\bfseries] at (5, 6.3) {\fa{۱۹۷۹-۸۰}};
\node[below, text width=2cm, align=center, text=DarkGray, font=\scriptsize] at (5, 5.4) {
    \fa{تاچر، ریگان}\\
    \fa{نئولیبرالیسم}
};

% فلش‌های غالب شدن
\draw[thickarrow, color=CenterGreen] (-1, 4) -- (1.5, 4);
\node[text=CenterGreen, font=\scriptsize] at (0.25, 3.6) {\fa{۱۹۴۵-۱۹۷۳}};

\draw[thickarrow, color=RightBlue] (2.5, 4) -- (6, 4);
\node[text=RightBlue, font=\scriptsize] at (4.25, 3.6) {\fa{۱۹۸۰-۲۰۰۸}};

% بحران ۲۰۰۸
\fill[OrangeAccent] (6.5, 6) circle (8pt);
\node[above, text=OrangeAccent, font=\small\bfseries] at (6.5, 6.3) {\fa{۲۰۰۸}};
\node[below, text width=1.5cm, align=center, text=DarkGray, font=\scriptsize] at (6.5, 5.4) {
    \fa{بحران}\\
    \fa{مالی}
};

\end{tikzpicture}
\caption{سیر تحول اندیشه اقتصادی از قرن هجدهم}
\label{fig:econ-history}
\end{figure}

\subsection{مکتب شیکاگو و نئولیبرالیسم}

\begin{personbox}[میلتون فریدمن (۱۹۱۲-۲۰۰۶)]

\textbf{ملیت:} آمریکایی

\textbf{اثر کلیدی:} \book{سرمایه‌داری و آزادی} (۱۹۶۲)

\textbf{جایزه نوبل اقتصاد:} ۱۹۷۶

\mediumspace

\person{فریدمن}، چهره اصلی مکتب شیکاگو، استدلال کرد:
\begin{itemize}[nosep]
    \item تورم همیشه و همه‌جا پدیده‌ای پولی است
    \item سیاست مالی کینزی ناکارآمد و حتی مضر است
    \item بازارها از دولت‌ها کارآمدترند
    \item باید خدمات دولتی را به حداقل رساند
\end{itemize}

\mediumspace

\textbf{نقل‌قول:} \textit{«اگر دولت مسئول صحرای صحارا شود، در عرض پنج سال کمبود شن خواهیم داشت.»}
\end{personbox}

\subsection{مکتب اتریش}

\begin{defbox}[تعریف: مکتب اتریش]
مکتب اقتصادی رادیکال‌تر از شیکاگو که ویژگی‌هایش عبارتند از:
\begin{itemize}[nosep]
    \item تأکید بر \textbf{فردگرایی روش‌شناختی}
    \item رد اقتصاد ریاضی و آماری
    \item تأکید بر \textbf{عدم قطعیت} و \textbf{کارآفرینی}
    \item مخالفت شدید با بانک مرکزی و پول فیات
    \item برخی به سمت \textbf{آنارکو-کاپیتالیسم} رفتند
\end{itemize}
\textbf{چهره‌های کلیدی:} میزس، هایک، راتبارد
\end{defbox}

\subsection{اقتصاد مارکسی معاصر}

اقتصاد مارکسی، با وجود فروپاشی شوروی، همچنان به حیات خود ادامه می‌دهد:

\begin{itemize}[nosep, rightmargin=1.5cm]
    \item \textbf{مکتب Monthly Review:} پل سوئیزی، تحلیل سرمایه‌داری انحصاری
    \item \textbf{مکتب تنظیم (فرانسه):} تحلیل بحران‌های ساختاری سرمایه‌داری
    \item \textbf{اقتصاد جهانی:} والرشتاین، نظام جهانی
    \item \textbf{اقتصاد فمینیستی مارکسی:} تحلیل کار خانگی و بازتولید
\end{itemize}

% ══════════════════════════════════════════════════════════════════════════════
\section{اختلافات درون‌مکتبی}
% ══════════════════════════════════════════════════════════════════════════════

مانند همیشه، اختلافات درونی هر جناح نیز مهم‌اند.

\begin{debatebox}[اختلاف درون چپ: بازار سوسیالیسم یا برنامه‌ریزی؟]
\textbf{سوسیالیست‌های بازار:}
\begin{itemize}[nosep]
    \item می‌توان مالکیت اجتماعی را با مکانیسم بازار ترکیب کرد
    \item تعاونی‌های کارگری در بازار رقابت کنند
    \item مدل یوگسلاوی، مدل مندراگون
\end{itemize}

\textbf{طرفداران برنامه‌ریزی:}
\begin{itemize}[nosep]
    \item بازار حتی با مالکیت اجتماعی، نابرابری و آنارشی ایجاد می‌کند
    \item فناوری امروز (رایانه، اینترنت) برنامه‌ریزی را ممکن‌تر کرده
    \item مدل برنامه‌ریزی مشارکتی (آلبرت و هانل)
\end{itemize}
\end{debatebox}

\begin{debatebox}[اختلاف درون راست: محافظه‌کاران و تجارت آزاد]
\textbf{لیبرتارین‌ها / نئولیبرال‌ها:}
\begin{itemize}[nosep]
    \item تجارت آزاد جهانی خوب است
    \item مهاجرت باید آزاد باشد (بازار کار جهانی)
    \item جهانی‌شدن رشد ایجاد می‌کند
\end{itemize}

\textbf{محافظه‌کاران ملی‌گرا:}
\begin{itemize}[nosep]
    \item تجارت آزاد، صنایع ملی را نابود کرده
    \item مهاجرت بی‌رویه، هویت ملی را تهدید می‌کند
    \item باید از کارگران داخلی حمایت کرد (تعرفه، محدودیت)
\end{itemize}
این اختلاف در عصر ترامپ بسیار پررنگ شد.
\end{debatebox}

% ══════════════════════════════════════════════════════════════════════════════
\section{جدول جامع: مکاتب اقتصادی و مواضع}
% ══════════════════════════════════════════════════════════════════════════════

% FIXED: Converted landscape table to portrait using tabularx for better integration
\begin{table}[htbp]
\centering
\caption{مقایسه مکاتب اقتصادی اصلی}
\label{tab:economic-schools}
\footnotesize % Reduced font size to fit 7 columns in portrait
\begin{tabularx}{\linewidth}{r|>{\centering\arraybackslash}X|>{\centering\arraybackslash}X|>{\centering\arraybackslash}X|>{\centering\arraybackslash}X|>{\centering\arraybackslash}X|c}
\toprule
\textbf{مکتب} & \textbf{بازار} & \textbf{دولت} & \textbf{مالکیت} & \textbf{پول} & \textbf{نابرابری} & \textbf{طیف} \\
\midrule
مارکسی & نقد & گذار & اشتراکی & زوال & محو & چپ \\
\midrule
پساکینزی & شکست جدی & بسیار فعال & مختلط & مدیریت & کاهش & چپ میانه \\
\midrule
کینزی & کارآمد & فعال & خصوصی & مدیریت & کنترل & میانه \\
\midrule
نئوکلاسیک & کارآمد & محدود & خصوصی & بانک مرکزی & نُرم & راست میانه \\
\midrule
شیکاگو & بسیار کارآمد & حداقل & خصوصی & قاعده‌مند & پذیرفتنی & راست \\
\midrule
اتریشی & تنها راه & ضرر & خصوصی & بازار & پذیرفتنی & لیبرتارین \\
\bottomrule
\end{tabularx}
\end{table}

% ══════════════════════════════════════════════════════════════════════════════
%                         خلاصه و پرسش‌ها
% ══════════════════════════════════════════════════════════════════════════════

\newpage

\begin{summarybox}

\textbf{۱. مالکیت:} بنیادی‌ترین اختلاف. راست از مالکیت خصوصی (حق طبیعی، کارایی) دفاع می‌کند؛ چپ مالکیت ابزار تولید را نقد می‌کند (پرودون: دزدی؛ مارکس: استثمار).

\textbf{۲. بازار vs برنامه‌ریزی:} هایک و میزس استدلال کردند بدون قیمت‌های بازار، محاسبه اقتصادی ناممکن است. چپ پاسخ‌هایی داده (سوسیالیسم بازار) اما تجربه شوروی، نقد راست را تقویت کرد.

\textbf{۳. نقش دولت:} از دولت حداقلی (لیبرتارین) تا دولت سوسیالیستی (مالکیت عمومی) طیفی وجود دارد. کینز با دولت رفاه، راه میانه‌ای پیشنهاد کرد.

\textbf{۴. نظریه ارزش:} مارکس ارزش را از کار می‌داند (پس سود = استثمار)؛ نئوکلاسیک‌ها ارزش را ذهنی می‌دانند (پس سود = پاداش).

\textbf{۵. مکاتب معاصر:} کینزگرایی پس از جنگ غالب بود؛ از ۱۹۸۰ نئولیبرالیسم؛ بحران ۲۰۰۸ بحث‌ها را بازگشود.

\textbf{۶. اختلافات درونی:} چپ بین بازار سوسیالیسم و برنامه‌ریزی اختلاف دارد؛ راست بین جهانی‌شدن و ملی‌گرایی اقتصادی.

\end{summarybox}

\vspace{20pt}

\begin{questionbox}

\begin{enumerate}[nosep, rightmargin=1cm]
    \item آیا می‌توان از مالکیت خصوصی دفاع کرد، اما همچنان ثروتِ موروثیِ بسیار بزرگ را ناعادلانه دانست؟

    \item آیا فروپاشی شوروی، پایان بحث «بازار یا برنامه‌ریزی» بود، یا شرایط دیگری (دموکراسی، فناوری) می‌تواند نتیجه متفاوتی دهد؟

    \item کینز خود را «لیبرال» می‌دانست، نه سوسیالیست. آیا دولت رفاه، میانه واقعی است یا چپ نرم؟

    \item نظریه ارزش کار و ارزش ذهنی، هر دو قابل دفاع به نظر می‌رسند. آیا این اختلاف قابل حل است یا صرفاً پیش‌فرض‌های مختلف؟
\end{enumerate}

\end{questionbox}

% ══════════════════════════════════════════════════════════════════════════════
%                         پایان فصل سوم
% ══════════════════════════════════════════════════════════════════════════════